% =====================================================================
%  GENERATED; DO NOT EDIT BY HAND.  Built by
%  python3 ebook-method/formel-logik/tr_parts.py build --write
%  from tr_template.tex and the two parts in ebook-method/formel-logik/tr/*.texfrag.
%
%  Harald Høffding, FORMAL LOGIC.  FOR USE IN LECTURES.  Fifth, revised edition.  English translation of
%  "Formel Logik. Til Brug ved Forelæsninger. Femte omarbejdede Udgave" (Copenhagen and Kristiania: Gyldendalske
%  Boghandel, Nordisk Forlag, 1907), made from the collated transcription (transcription.tex in this folder) by
%  two translator agents (2026-10-29) working from GLOSSARY.md and TRANSLATOR-BRIEF.md (ebook-method/formel-logik/);
%  their 19 odd spots were settled at the image and the Danish mended where it was the transcription's fault.  The printed page numbers
%  of the 1907 edition are in the margin.  The logical notation (math mode, \sim, \udel, the schemes) is the
%  transcription's, unaltered; tr_parts.py build checks it character for character.  Danish example sentences are
%  translated, with the Danish word in brackets where the logical point turns on the Danish wording.
% =====================================================================
\documentclass[12pt,a4paper,openany]{book}
\usepackage[utf8]{inputenc}
\usepackage[T1]{fontenc}
\usepackage{libertinus}
\usepackage{libertinust1math}
\usepackage{array}
\usepackage{arydshln}   % formel-logik: \cdashline, the short dotted rule of the schemes pp. 44-45 (notatex.py)
\usepackage{graphicx}
\usepackage{newunicodechar}
\newunicodechar{ɔ}{\reflectbox{c}}   % formel-logik: »ɔ:« (det er), as in the Brøchner transcriptions
\newcommand{\udel}{\mathbin{\cdot\kern-0.1em|\kern-0.1em\cdot}}   % formel-logik: the exclusion sign ·|· of p. 27 (»gensidigt Udelukkelsesforhold«)
\usepackage{textalpha}
\usepackage[english]{babel}
\usepackage[protrusion=true,expansion=true]{microtype}
\usepackage{setspace}
\usepackage[margin=1.3in]{geometry}
\usepackage{fancyhdr}
\usepackage{hyperref}
\hypersetup{pdftitle={Formal Logic}, pdfauthor={Harald Høffding}, hidelinks}
\pagestyle{fancy}
\fancyhf{}
\fancyhead[LE,RO]{\thepage}
\fancyhead[C]{\textit{Høffding: Formal Logic, 5th ed. (1907)}}
\setlength{\headheight}{15pt}
\setstretch{1.3}
\usepackage{marginnote}
\newcommand{\opage}[1]{\marginnote{\footnotesize\textit{[#1]}}}
\newcommand{\whiteline}{\par\vspace{\baselineskip}\par}
% heading macros (from ledende-tanker; to be used by patch once the headings are read at the image)
\newcommand{\parthead}[1]{\par\bigskip\phantomsection\begin{center}{\Large #1}\end{center}\addcontentsline{toc}{part}{#1}\medskip}
\newcommand{\unnumhead}[1]{\par\bigskip\phantomsection\begin{center}{\Large #1}\end{center}\addcontentsline{toc}{chapter}{#1}\medskip}
\newcommand{\chapterhead}[2]{\par\bigskip\phantomsection\begin{center}{\large #1\ #2}\end{center}\addcontentsline{toc}{chapter}{#1 #2}\medskip}
\newcommand{\chapterheadii}[3]{\par\bigskip\phantomsection\begin{center}{\large #1\ #2\\ #3}\end{center}\addcontentsline{toc}{chapter}{#1 #2 #3}\medskip}
\newcommand{\sectionhead}[2]{\par\bigskip\phantomsection\begin{center}#1 \textit{#2}\end{center}\addcontentsline{toc}{section}{#1 #2}\medskip\ignorespaces}
\newcommand{\ix}[3]{\par\noindent\hspace*{#1}#2\dotfill #3}
% formel-logik (heads.py): chapter = numeral, title in capitals, the smaller subtitle; section heads; letterspacing
\newcommand{\chaphead}[3]{\par\bigskip\phantomsection\begin{center}\ifx&#1&\else{\large #1}\par\smallskip\fi{\Large #2}\ifx&#3&\else\par\smallskip{\large #3}\fi\end{center}\addcontentsline{toc}{chapter}{\ifx&#1&\else #1. \fi #2}\medskip}
\newcommand{\sechead}[2]{\par\bigskip\phantomsection\begin{center}#2\end{center}\addcontentsline{toc}{section}{#1}\medskip}
\newcommand{\spaced}[1]{\textls[180]{#1}}

% the translator's note, set by tr_parts.py build before the Forord (p. VII)
\newcommand{\translatorsnote}{\chapter*{Translator's Note}
\noindent This is a translation of the fifth, revised edition (1907) of Harald Høffding's \textit{Formel Logik}.  The
printed page numbers of that edition are in the margin.

Høffding's logical notation is kept exactly as printed.  Capital letters ($A$, $B$, $C$, \dots) stand for concepts, and
the corresponding small letters for their negations ($a$ = non $A$).  Letters written together ($AB$) stand for the
combination of concepts, the concept whose intension is the sum of theirs.  The sign $=$ means identity, that is, the
possibility of substitution ($A = AB$: $A$ does not occur without $B$); $<$ and $>$ compare extensions ($A < B$: the
extension of $A$ is part of that of $B$); $\sim$ (Danish ``er ikke,'' ``is not'') marks the contrary relation; $\udel$
the relation of mutual exclusion; $0$ the self-contradictory combination, ``nothing.''  Raised numerals
(\textsuperscript{1}), \textsuperscript{2}), \dots) number the cases of an inference.

Danish example sentences are translated; where the logical point turns on the Danish wording, the Danish word is given in
square brackets.  Misprints of the 1907 printing are translated in their evident sense.
\clearpage}

\begin{document}

\frontmatter
% --- p. III ---
\thispagestyle{empty}\opage{III}\vspace*{5cm}
\begin{center}{\Large FORMAL LOGIC}\end{center}
\clearpage

% --- p. V ---
\begin{titlepage}\centering\opage{V}
\vspace*{1cm}
{\large HARALD HØFFDING}\par\medskip
\rule{3cm}{0.4pt}\par\vspace{1.5cm}
{\Huge FORMAL LOGIC}\par\bigskip
{\large FOR USE IN LECTURES}\par\vspace{1.2cm}
\textit{FIFTH, REVISED EDITION}\par
\vspace{2cm}
\vspace{3cm}
COPENHAGEN AND KRISTIANIA\par
GYLDENDAL BOOKSHOP\par
NORDISK FORLAG\par\medskip
\rule{1.5cm}{0.4pt}\par\medskip
{\small THE PUBLISHER'S PRINTING HOUSE}\par\medskip
1907\par
\vfill
{\small Translated from the Danish}\par
\end{titlepage}
\translatorsnote
% --- p. VII ---
\opage{VII}

\chapter*{PREFACE TO THE FIFTH EDITION}


In its manner of presentation this outline of logic follows
the logic of identity and of intension set forth by \textsc{Leibniz}, \textsc{Boole} and \textsc{Jevons}.
Through repeated occupation with this manner of presentation
and its application in my teaching,
I have become confirmed in the view that it aims more closely
at the essential side of our thinking, and that
by its simplicity and clarity it is better suited to pedagogical
use than the usual manner of presentation. ---

In the present fifth edition I have undertaken a thoroughgoing
revision, based on the views I
have developed in my treatise \textit{Det psykologiske Grundlag
for logiske Domme} (Videnskabernes Selskabs Skrifter. 6th
series) and in the essay on the concept of analogy (Mind
1905, in Danish in the second series of my ``Mindre Arbejder'').
I have also been able to make use of some valuable
critical remarks from one of my former listeners
and from Dr.\ \textsc{Anton Thomsen}.

\textit{December 1906.}

\hfill\textbf{Harald Høffding.}

% --- p. IX ---
\chapter*{\opage{IX}CONTENTS}
\begin{small}


\ix{0em}{\textbf{I. Synthesis and Analysis.}}{1}
\par\noindent\textbf{II. Analysis of Intuitions:} Formation of Concepts and Judgments.
\ix{1em}{A. Formation of Concepts}{5}
\ix{1em}{B. Intension and Extension of Concepts.}{8}
\ix{1em}{C. Formation and Formulation of Judgments.}{11}
\ix{1em}{D. Secondary Judgments.}{18}
\par\noindent\textbf{III. Analysis of Judgments:} Inference.
\ix{1em}{A. The Logical Principles}{23}
\ix{1em}{B. Inferences}{28}
\ix{2em}{a. Direct Inference}{29}
\ix{2em}{b. Indirect Inference.}{34}
\ix{3em}{α. Indirect Inference by Pure Deduction}{35}
\ix{3em}{β. Indirect Inference as a Union of Induction and Deduction}{38}
\par\noindent\textbf{Appendix:}
\ix{1em}{1. Mill's Methods of Induction}{42}
\ix{1em}{2. Analogy and Inference by Analogy}{46}
\end{small}
\clearpage

\mainmatter

% --- p. 1 ---
\opage{1}\chaphead{I}{SYNTHESIS AND ANALYSIS}{}


1. To think is to compare, to find difference and similarity.
In different forms and to different degrees such
a comparing shows itself in all perceiving and having of ideas, in all
observation and connection of ideas. Indeed, every state of consciousness
is conditioned by the differences and similarities found
within it and in its relation to other states of consciousness.
Pleasure and pain are conditioned by an experience's relation to our nature
and our condition; to will is to prefer, and so likewise rests
on a relation of difference. In a narrower sense only the
comparison carried out with explicit, concentrated (involuntary or voluntary) attention
is called thinking. Such
attention is awakened when the task is to solve a problem
or to find means to an end. (Cf.\ Psykologi VB, 11.)

Comparison presupposes energy to be able, at one stroke or in
immediate succession, to hold fast to and hold together
different elements; only through such holding together do
the similarity or the difference appear. A psychical energy that manifests
itself as synthesis is therefore the ultimate presupposition of all thinking.
(Cf.\ Psykologi II, 5.)

Seen from another side too, synthesis is the ultimate
presupposition of thinking. Comparison takes place by attention
passing from element to element; but this presupposes that the
different elements, prior to the comparison, are given
in connection, in a total intuition. Only on the basis of such
a total intuition can analysis, the particular turning
of attention to the individual elements (parts, sides
% --- p. 2 ---
\opage{2}or qualities), take place. --- The total intuition may be
sense-intuition, memory-intuition or imagination-intuition. Thinking
can take place on the basis of all these different
kinds of intuition, and it then becomes a special question
how these are related to one another. --- In every such
total intuition a manifold of elements is connected.


The comparing and connecting activity which, seen from two sides in this way,
is the condition of comparison, and hence of
thinking, we call synthesis. \textit{All analysis presupposes synthesis,
just as all comparison does}.

Total intuitions are formed from the first involuntarily, and we hold
fast to them as long as possible. Only reluctantly do we proceed
to dissolve them. But another ordering of the elements
than the one given in the involuntary intuitions may be
necessary, partly in order to give us an overview of the many
different intuitions we have, partly in order to remove the contradictions
that may appear within the same intuition or
in the relation between several different intuitions which occur
simultaneously or close at hand.

Whereas psychology, as the natural history of thinking, investigates
the actual thought processes that take place in us, logic seeks
rules and standards by which a distinction can be made
between valid and invalid thinking. Valid thinking must
be that which arrives at firm and universally valid results. By
a firm result is meant one that can be held fast without
the thinker coming into contradiction with himself (even if
ever so many sense-, memory- and imagination-intuitions crop
up); by a \textit{universally valid} result is meant one that is
recognized not merely by the thinker himself, but by other
thinking beings at the same stage of development as he. Logical
thinking springs from a striving to preserve
identity with ourselves and with all those whose power of thinking we
recognize as equal to our own. \textit{Valid thinking is
that by which it is possible to maintain as great a correspondence
as possible among as many elements of consciousness
(our own and others') as possible}. Experience shows that the more we
% --- p. 3 ---
\opage{3}hold fast to the standard thereby given, the fewer disappointments we suffer,
and the more secure the expectations in which we can indulge. ---

Logical thinking is in a certain sense not productive.
Its work consists in bringing to explicit consciousness, through analysis and comparison,
what is already contained
in the given total intuitions, which are the work of the original synthesis.
It brings the content given in our sense-, memory- and
imagination-intuitions into new connections, called
concepts and judgments. From judgments inferences can in turn be drawn;
in which the relation between the concepts contained in different judgments
comes to explicit consciousness; but a wholly new
content the inference does not bring in relation to the judgments, any more
than judgments and concepts do in relation to the intuition. And
yet it can be not merely a strenuous, but also a significant
psychical work that has been done when the content of consciousness
has been recast into new forms (from intuitions to concepts and
judgments, and from judgments to inferences), so that new syntheses
can take the place of the first syntheses. Thereby
not only can an overview and clarity be won of what has hitherto been intuited,
but a need and capacity can also be awakened to seek out
and produce new intuitions.

2. \textit{Every} science is a doctrine of valid thinking in
its own domain. Formal logic investigates the general
forms of valid thinking, the conditions of thought's agreement
with itself, regardless of which domains thinking
moves in. It abstracts from all determinate content
and can therefore take its examples from all possible domains.
It is, after all, always one and the same human thought that is at work,
and there must then be certain rules of validity that hold everywhere.

Mathematics too may be called a formal science, since it
abstracts from the specific qualities of objects and merely investigates
similarities and differences between figures and magnitudes.
But logic is formal to a still higher degree, since even figures
and magnitudes are for it specific qualities; it investigates
similarity and difference in all generality, wherever they appear
to thinking. The logical principles therefore form
presuppositions for mathematics as for every other science.


% --- p. 4 ---
\opage{4}Formal logic does not coincide with the theory of knowledge,
which is also sometimes called logic. The former tests the agreement
among our ideas without asking about
their real significance. Logical thinking can be practiced on the basis
of imagination-intuitions as well as on the basis of memory- and
sense-intuitions. Formal logic passes over into the theory of knowledge
when the question is asked whether the forms
and results of our thinking hold for real existence: whether there is consistency
in existence, as there is consistency in what logic
regards as correct thinking. This question arises
naturally from the fact that thinking is practiced on the basis of intuitions
which it has not itself produced. In connection with this
there also arises the question of the relation among sense-,
memory- and imagination-intuitions, a relation which formal logic
leaves aside.

% --- p. 5 ---
\opage{5}\chaphead{II}{ANALYSIS OF INTUITIONS}{FORMATION OF CONCEPTS AND JUDGMENTS}


\sechead{A. Formation of Concepts.}{\textbf{A. Formation of Concepts.}}


3. The first condition of valid thinking is that we become
distinctly conscious of what is contained in the intuitions
formed by involuntary synthesis. This is attained by analysis, that is,
by a series of acts of attention (see 1). Of the table
at which I sit I have an image, which by an involuntary
synthesis is formed from immediate and reproduced sensations
of sight, touch and movement. I analyze it when
I notice that it is black, hard, square, long, etc. A
chemical substance (e.g., carbonic acid) I determine by taking note of
its appearance, its effects, its occurrence in nature,
its origin. In psychology I characterize consciousness
by taking note of its limiting states and of its peculiar
mode of operation (see Psykologi II, 5).

4. This analysis is supported partly by comparison among
the elements (the parts or qualities) themselves, partly by
comparison between them and similar or different elements
in other intuitions, which are reproduced (partly in bound,
partly in free form, see Psykologi VB, 1---2). Only thereby can the
distinctive character of the individual element properly appear. The location and magnitude
of the individual part can be determined only when it
is held together with other parts; the individual quality must be
held together with others in order to become distinct. (Cf.
the law of relation. Psykologi VA, 5; D, 5). Often an element
can be distinguished only when it has occurred in another connection,
% --- p. 6 ---
\opage{6}so that a recognition can take place. The color and shape of the table
I notice rather when I have seen the same color with another
shape and the same shape with another color. This recognition
is expressed by my giving the qualities names. If I call the
table square, it means that I recognize a shape I
have met with in other intuition-images. When Newton called the
force that keeps the planets in the vicinity of the sun gravitation,
he wished thereby merely to signify that it is the same
force that brings the hurled stone back to the earth.

Analysis thus rests on a synthesis, which works as
an auxiliary process. Since intuition and comparison already presuppose
synthesis (1), we see how analysis and synthesis stand in
constant interaction.

5. Having now found the elements contained in the
given intuition, the task is to connect and order them in such a way
that we obtain a unified characterization of the given content. There
must then again be performed a gathering together, a synthesis; only
thereby is the whole process concluded which leads to the framing of \textit{a
concept}.

From the mere idea, which is only a reproduction of a
given intuition (in the simplest case only of a sensation),
the concept differs by the clarity and distinctness won through the interaction of
analysis and synthesis. A concept
is an idea whose content we have made clear and distinct to ourselves,
so that it is not altered by the different connection
in which it occurs. Our consciousness can now at
this point preserve its identity with itself, and accidental
associations of ideas can be prevented from intervening. Thereby
it also becomes possible for several individuals to be able to think
together: only when each for himself has fixed his ideas
can they connect the same ideas with the same word.


\textit{A concept is thus an idea that is firm (for the same individual) and
universally valid (for the different thinking individuals)}.
--- The definition (the fixing of the concept, the determination of the concept)
consists essentially in a delimitation of the content of ideas.
It seeks to reduce this to the least possible by
% --- p. 7 ---
\opage{7}omitting the elements that are not necessary to distinguish the one
content of an idea from the other.

A concept may be formed on the basis of a single
intuition, or on the basis of several different intuitions
of the same phenomenon (thing, person, event, relation), or
on the basis of several intuitions that have different phenomena
as their object. In the first two cases the concept is called
an \textit{individual concept}, in the last case a general concept.
(Cf.\ Psykologi VB, 9.) If the individual concept is formed on the basis
of a single intuition, it is called \textit{concrete}; if it is formed
on the basis of several different intuitions, it is called \textit{typical}.
General concepts may be of lower or higher degree, according
as they stand nearer to or farther from immediate intuition;
the former are called \textit{species concepts}, the latter \textit{generic concepts}. What in one
respect is species may in another respect be genus, and
vice versa.

The \textit{extension} of a concept is the sum of the intuitions for which
it holds; the intension of a concept is the marks (qualities
or parts) which have been found by analysis of the given
intuitions.

Where several different intuitions underlie it, that is,
in typical and general concepts, the similarity that
is found among the individual intuition-images need only be similarity of relations
(analogy): it will be the relation among the parts
or qualities that is similar, not the individual parts or
qualities themselves. In the individual concept of a particular human being,
which is to hold for him as young, as grown
and as old, the similarity on which the concept builds is similarity of relations.
Likewise in the concept horse, the concept vertebrate, etc.
The concept expresses the law of the mutual relation of the parts or qualities.
Only where the concept indicates a quite simple
mark (as e.g. the label on trunks that are registered
from Copenhagen to London) can similarity of coincidence or similarity of quality
be present. Where this is not the case, it is only
by the aid of analogy that consciousness can go from example to
example while preserving its identity with itself. (Cf.
on similarity of coincidence, of quality and of relations Psykologi VB, 8.)
% --- p. 8 ---
\opage{8}The simplest and clearest example is taken as the basis, perhaps
merely in schematic form, i.e., with omission of all dispensable
elements. The concept is a direction for carrying out
a movement of thought through a series of intuitions. Whether
it can be completed depends on whether the concept is formed on the basis
of sufficiently many and sufficiently complete intuitions.
Where the content of the intuition is inexhaustible, the definition
becomes merely a hypothesis, which may be necessary for a provisional overview,
but must find its confirmation through continued analysis.

6. The formation of concepts takes place by the elements of a given intuition
being seen in their mutual relation and in relation to the elements
in other intuitions, and the determination of the concept takes place
by a delimitation of the content found by analysis.
\textit{Every concept therefore expresses a relation}. I become
conscious of the elements in the content of an intuition by considering
them in their difference and similarity, partly over against one another,
partly over against elements in other intuitions. Every
concept must therefore be defined in relation to other concepts.
A definition must be changed, partly when the content of the intuition
changes, partly when the nearest lying concepts change.
There is therefore a constant interaction between concept and intuition
and between definition and classification. The concept of chemical
composition changed when Boyle demonstrated how, by
the composition of elementary substances, a substance with
other properties than these can arise; and thereby it became further
possible to distinguish sharply between chemical connection
and mechanical mixture. Rousseau's demonstration of the distinctive character
of the life of feeling led to a new division of psychology.

If one cannot point out the content of a concept that forms
a limit for a given concept, one forms provisionally a negative
concept (non $A$ as limit for $A$). See below: 8.

\sechead{B. Intension and Extension of Concepts.}{\textbf{B. Intension and Extension of Concepts.}}


7. \textit{The intension and extension of concepts} (5) \textit{stand in inverse relation
to one another}. For the more different intuitions that underlie
the formation of concepts, the fewer marks can
% --- p. 9 ---
\opage{9}be included in the intension of the concept. The more different intuitions
a concept is to hold for, the fewer qualities it can assert
of them. Thus the typical individual concept has greater
extension but less intension than a corresponding concrete individual concept;
a general concept has greater extension but less
intension than an individual concept, and in the same way it goes
with species concepts and generic concepts in relation to one another. A concept
is called lower in relation to another when it has less extension
but greater intension, --- higher when it has greater extension
but less intension. Thus the concept human being
is higher in relation to the concept Negro: it comprises more individuals
than the latter, but it cannot specify as many marks
as the latter.

The proposition about the inverse relation between extension and
intension holds only when we compare concepts that lie
in the same series, i.e., when they find application within one
and the same group of intuitions and marks, without however
being coordinate. It does not hold for two coordinate species concepts;
for the extension of a species can be widened or narrowed
without its intension being diminished or increased. Because there are
more sparrows than eagles in the world, the concept sparrow
need not have less intension than the concept eagle. The concept
bird, on the other hand, has greater extension but less intension than
both of the two concepts mentioned. Other examples of concepts
that lie in the same series are: Negro --- human being --- living
being; hatred --- feeling --- element of consciousness; square --- rectangle
--- quadrilateral --- rectilinear figure --- figure.

It must also be remarked that we do not always think a concept
with all its marks, but only with those that are necessary
to distinguish it from the nearest lying concepts.
We content ourselves with defining the human being as a rational being
when it is merely a question of distinguishing the human being from the animal. The law
of the inverse relation between intension and extension holds
of course only when we take the whole intension into account.

8. \textit{Negative concepts} are formed when, in the attempt at
definition, it turns out that there is a limit to the intension and extension
of a concept, without our being able for the time being to specify positive
% --- p. 10 ---
\opage{10}marks for that which lies on the other side of the limit.
In classification we must then content ourselves with entering
a concept with negating marks: non $A$ as limit
for $A$. (This notation recalls the concepts mentioned in 5
whose only intension is a simple mark.) The history of classifications
shows that one often has to form groups for the time being without
positive marks and leave it to later research to
find such, which will then lead to a further division of
the manifold that is perhaps comprised in the negative concept.
Aristotle divided the animals into rational and irrational,
and the irrational again into those that have (red) blood and those that
do not have (red) blood. Often a negative concept may be concealed
under a positive designation; thus Linnaeus, for his
sixth class of animals, which he called the worms, could not specify other
marks than that they were neither mammals, birds, amphibians,
fishes nor insects. Later Lamarck divided the animals into
vertebrates and invertebrates, and Cuvier was the first to find positive
marks for and a positive division of the invertebrates.
--- Negative concepts can thus have the significance of marking
provisional limits in the development of research.

The relation between the intensions of two negative concepts is the
converse of the relation between the intensions of the two corresponding positive concepts.
The concept ``living being'' has less intension
than the concept ``human being,'' but the concept ``non-living'' has
greater intension than the concept ``non-human.'' Thus $A < B$,
but $a > b$ (when the small letters denote the negative concepts).
--- The corresponding holds for the extension. The concept
``living being'' has greater extension than the concept ``human being,''
but the concept ``non-living'' has less extension than the concept
``non-human.'' $A > B$, but $a < b$. --- The negations of the higher
concepts have greater intension but less extension than the negations of the
lower concepts.

9. In the use of the words ``higher'' and ``lower'' regard is had
to the extension of concepts, and older logic has
on the whole laid predominant weight on this side of the concept.
And yet it is the intension that first and foremost plays a role
% --- p. 11 ---
\opage{11}in thinking. It must first be decided which marks
are to be included in the concept, before it can be decided in how
many different cases these marks are found again.
This is connected with the fact that we form our concepts from the outset
on the basis of intuition; when we proceed strictly methodically,
we must therefore begin with concrete individual concepts.
It is from the outset undetermined how great the extension
may become, --- whether we can form typical and general concepts.
It is thus the intension we begin with when we proceed logically.
--- Psychology, to be sure, teaches us that the art may be precisely
to grasp the concretely individual in its whole distinctiveness,
and that a vague grasp of the particular misleads
into hasty generalization, so that the development our ideas
must undergo, seen from one essential side, consists in
an individualization and specialization. (Psykologi VB, 9b.) But
this shows only that the logical doctrine of valid thinking is
different from the psychological doctrine of actual thinking.

\sechead{C. Formation and Formulation of the Judgment.}{\textbf{C. Formation and Formulation of the Judgment.}}


10. Formation of judgments and formation of concepts cannot be separated
from one another, but are two sides of the same thought process. The
judgment too is formed, where it is formed from the outset, by analysis
of intuition. Whereas in the formation of concepts the analysis aims
at emphasizing and isolating individual elements, which come
to make up the intension of the concept, in the formation of judgments it aims
at determining the relation among the elements given in the intuition.
But to a distinct concept there also belongs the determination
of this relation, and the very framing, the definition
of the concept, takes place by a judgment, since the individual concept can only
be determined by its relation to other concepts. (Cf. the relation
between definition and classification: 6. 8.) The auxiliary
synthesis (4) plays a greater role in the formation of judgments than
in the formation of concepts, since the grasp of the mutual relation
among the elements of the intuition, as well as of the relation of the intuition
to other intuitions, presupposes a more pronounced holding together
with something that is not immediately given.

% --- p. 12 ---
\opage{12}A judgment is a \textit{firm and universally valid connection of concepts}.
(Cf.\ 5.)

According to the legend, Newton, at the sight of a falling apple,
is supposed to have hit upon the thought that gravity must operate throughout the whole
of space, hence on the moon in relation to the earth and on
the planets in relation to the sun, and from the intuition-image thus arisen
he drew the judgment which he later
sought to prove: planetary motion is a motion of falling. --- If
I say: this man is good! --- this judgment is derived (when I
pronounce it on first impression) from my sense- and memory-intuitions
of the man and his conduct, in which I find the quality
``good'' (which I know from other intuitions, see 4).

For a proof of the validity of a judgment one must go the way of
inference, through inductions and deductions. The relation
between intuition and judgment is, to be sure, analogous to the relation
between ground and consequence, but is not so fully and clearly
conscious as the latter can be.

This analogy also holds in the respect that, just as
the consequence contains only what was already contained in
the ground, so the judgment brings no wholly new content,
but only a new, more determinate and conscious form for the
content given in the intuition.

11. The relation which the judgment asserts as holding between
two concepts (or elements of concepts) can be expressed in different
ways. One can keep to \textit{the usual linguistic
mode of expression}, where the one concept or element (the predicate)
is referred to the other (the subject) by means of the word
``er'' [is] (copula). Or one can keep to this, that the one concept
(the predicate concept) has a greater extension than the other (the subject concept),
so that the latter is included in (subsumed under)
the former. This \textit{relation of subsumption} can then be illustrated
by geometrical figures, so that e.g.\ the subject is represented as
a smaller circle inside a larger one, which denotes the predicate.
(An analogous mode of presentation could be applied with regard to
the relation between the intensions of concepts, so that the predicate concept
were expressed as a small circle inside the larger
circle of the subject concept. This mode of presentation has however never been tried.)
% --- p. 13 ---
\opage{13}Or one can formulate the judgment as the expression of \textit{a relation of identity}
between two concepts, hence as an equation. This
mode of expression recalls algebra, just as the former recalls
geometry.

But although the judgment can recall grammar,
geometry and algebra alike, it is nonetheless in itself independent of these
modes of expression, as is seen from the fact that words, figures or signs
must be chosen according to the relation of thought and determined by it. A
closer demonstration of this will serve to illuminate the relation of thought
to the manner in which it is expressed.

12. The logical subject is the more indeterminate concept
from which the movement of thought in the formation of judgments proceeds,
and about which it turns (the initial idea); the logical predicate
is the concept in which the movement of thought ends (the final idea),
because the subject has found in it its closer
determination. The Danish language connects them by the word ``er'' [is].
But the grammatical sentence is not always a simple formulation
of such a relation of concepts.

a) \textit{A grammatical sentence can contain several logical judgments
at once}. The sentence ``This sofa is hard'' can
contain three judgments, since the statement can contrast this sofa
with another sofa, or characterize the object pointed at as a sofa,
or merely assert hardness of it.
``Gustavus Adolphus fell at Lützen on November 6, 1632'' asserts
not merely the King's fall, but time and place. b) This
plurality of judgments in one and the same sentence is connected
with the fact that \textit{the grammatical predicate is not always the logical
predicate}. One can find the logical predicate by
noting where the stress lies. If one says ``\textit{this} sofa is hard,''
the logical predicate lies in ``this,'' and the judgment becomes: ``the
sofa that is hard is this one here.'' ``\textit{All} the guests have come'':
here the logical predicate lies in ``All''; it is the new element. ``\textit{Then}
the powder magazine blew up'': the logical predicate here is the point of time:
it was known beforehand that the powder magazine had blown up.
In sentences that are not uttered aloud, and where the stress
is not indicated, interpretation from the whole context must find
the logical predicate. c) \textit{There are sentences in which the logical}
% --- p. 14 ---
\opage{14}\textit{subject is not expressed at all}, because it goes without saying
or lies in the situation. Such predicate judgments
we have in the simple exclamations like ``Splendid!'' ``Bravo!'' or in
announcements like ``The King!'' --- To this class belong also the indeterminately
descriptive sentences in which merely a state or an activity
is stated: ``There's a draft.'' ``It is declining.'' ``There is talk
of ....'' The initial idea is here indeterminate and cannot be
formulated, or it is self-evident, and the interest is wholly absorbed
in the logical predicate.

The significance of stress and the presence of predicate judgments
show that the psychological interest is attached to the
logical predicate, which brings what is new, whereas the logical
subject stands as provisionally known, even though it needs
closer determination. If a clear consciousness is to be reached of
the relation the judgment asserts, this relation must however be formulated in such a way
that both subject and predicate are determined expressly;
only then can the relation between them become clear, the relation
to which the logical interest is properly attached. Only when
this relation is distinctly stated can the judgment really be
a firm and universally valid connection of concepts. Logically viewed, it is
therefore in the end without significance what is the initial idea
and what is the final idea, and how we
get from the former to the latter. The temporal relation loses its significance.
The grammatical formulation stands in this respect nearer
to psychology than to logic. The two following formulations
have, on the other hand, a purely logical character.

13. When one, like the older logic from Aristotle on,
takes the extension of concepts as the basis, the judgment means that
the group of phenomena for which the one concept (the subject concept)
holds is a part of the group of phenomena
for which the other concept (the predicate concept) holds. The judgment thus expresses
\textit{a relation of subordination (relation of subsumption)}.
That Negroes are human beings means that the group of Negroes is smaller
than the group of human beings and is included in it. One can illustrate
this relation by representing the higher concept as
a larger circle, within whose circumference the other concept is represented
as a smaller circle. That $A$ is $B$ thus means $A < B$.

% --- p. 15 ---
\opage{15}In this way one will be able to give a valid presentation
of logic; only such judgments in which the two concepts
are absolutely identical ($A = B$) will one be unable to find room
for, and they do not occur in Aristotle either. Geometrically
we should here have only a single circle and could not see that
it means the coincidence of two concepts. ---

Because one starts from the extension of concepts, one need nevertheless not
conceive the judgment as subsumption. For $A < B$ implies
that $A$ is identical with a part of $B$'s extension. Negroes are a part
of humanity. Let $n$ now mean ``some,'' then $A = nB$. For
$n = 1$ (where 1 means the whole totality as opposed to a
part or some) we get $A = B$. --- Thereby the judgment is represented
as the expression of a relation of identity between two extensions.

But it is then natural to go a step further. When $A$ is
$B$, $A$'s extension is perhaps only a part of $B$'s extension; but
in any case $B$'s whole intension must hold for $A$. And since it
is the intension of the concept that thinking first turns to,
while the determination of the extension is secondary (9), it will be
natural to take the intension of concepts as the basis. This leads
us to the third mode of presentation, which expresses the judgment
by a relation of identity between intensions of concepts.

14. Just as the need to form concepts springs from
consciousness's need to preserve its agreement, if
possible its identity with itself, despite the manifoldness and
change of the content of consciousness, so there will also be a natural
need to preserve identity when we connect one concept with
another in a judgment. The more the two concepts stand external and
foreign to one another, the less will consciousness be able to
preserve its unity, --- the more will the connection stand as forced,
in the end as a mere verbal connection. The more, on the other hand, we
can see the relation between the two concepts as a
relation of identity, the more natural and clear will the connection
be. When we determine more closely the different degrees or
kinds of identity that may be present, we shall be able to
maintain the contention that \textit{the judgment expresses a relation of identity}.
--- Leibniz and (in the 19th century) Boole and Jevons have
founded this mode of presentation.

% --- p. 16 ---
\opage{16}There are three different relations of identity.

a) The intension of the predicate concept may be a closer determination
of the intension of the subject concept, without the two concepts
coinciding entirely. When I, e.g., by observing a man's
conduct, am led to pronounce the judgment ``The man is good,'' the concept
``good'' is a closer determination of the individual concept
``the man.'' The goodness is found in the man's conduct, and only because
it is there is the judgment asserted. In the designation ``good'' lies a
reference (by auxiliary synthesis, see 4 and 10) to other phenomena
which have presented a similar quality. By the judgment
``The orbit of Mars is an ellipse'' Kepler, after an analysis
of Tycho Brahe's observations, pronounced a more determinate characterization
of the planet's orbit than had hitherto been had; by the designation
``ellipse'' there was expressed a recognition (by means of auxiliary
synthesis) of a figure that was known from the doctrine of
conic sections. --- This kind of judgment can be formulated thus:
$A = AB$. The sign of equality here means possibility of substitution,
i.e., of putting the one expression in place of the other. The juxtaposition
of the letters ($AB$) means a concept whose
intension is the sum of the intensions of two concepts ($A$'s and $B$'s),
hence a combination of concepts. $A = AB$ thus means that when
we have $A$, we have a richer intension than we thought. $A$ does not occur
without $B$, but $B$ may occur without $A$. There
are other ellipses besides the orbit of Mars, and other good beings
besides that man. $B$ has a greater extension than $A$: $A < B$.
--- Perhaps $A$ has more determinations than $B$, and we then have
$A = ABCD$ .....

A relation of identity of this kind is called \textit{partial identity}.

b) When not only is $B$ a closer determination of $A$,
but $A$ also proves to be a closer determination of
$B$, so that we have both $A = AB$ and $B = BA$, there is \textit{absolute
identity} between $A$ and $B$. It is immaterial in what order
the concepts occur in a combination, since this does not alter
the sum of the intension. Therefore $AB = BA$.

Two absolutely identical concepts are only equally valid expressions
for one and the same intension of a concept, which can be considered from two
different sides. Examples: ``This house is my home.'' ``It
% --- p. 17 ---
\opage{17}is Peter who is coming'' (the one coming = Peter). The King of
England = the Emperor of India. Grundtvig's year of birth = the
fourth year after Oehlenschläger's year of birth. The circle = the
line returning into itself, in which every point has
the same distance from one and the same point.

Whereas the demonstration of partial identity presupposes only an
analysis of the one concept ($A$), since the other ($B$) is contained
in it, the demonstration of absolute identity presupposes an analysis
of both concepts. The thought activity by which absolute
identity is discovered is thus more composite than that by
which partial identity is discovered. An absolute identity presupposes
two partial identities ($A = AB$ and $B = BA$). Nevertheless
there is an involuntary tendency to conceive a relation of identity
as absolute before it has been tested from both sides.

c) The relation between two concepts ($A$ and $B$) may finally
be that $A$ is partially or absolutely identical with $B$
only under certain conditions. The proposition ``Water is liquid,'' holds
only when the temperature is above 0°, but not above 80° (Réaumur).
--- This relation is called \textit{limited identity}. It is expressed
thus: $AC = ABC$ (limited partial identity) or
$AC = BC$ (limited absolute identity). $C$ here denotes the
condition without which the relation of identity does not hold.

On closer investigation it will turn out that every judgment
expresses limited identity (just as every concept expresses
a relation, see 6). We omit the limitation where it does not
appear distinctly, and where it is not of essential significance
for the course of thought. Often special experiences are needed in order that
the limitation be discovered. An inhabitant of the tropics would not think
that the proposition ``Water is liquid'' needed any limitation
(provided he does not call steam water).

The judgments that assert limited identity presuppose not merely
an analysis of one or both concepts, but also of the whole
larger context within which their intension occurs,
and within which the conditions for their
identity must also be found.

% --- p. 18 ---
\opage{18}\sechead{D. Secondary Judgments.}{\textbf{D. Secondary Judgments.}}


15. When on the basis of intuition a judgment has been formed,
the question can be raised of the validity of this judgment, --- whether
it holds for all phenomena of the same kind, --- whether it
holds unconditionally, --- whether it holds at all, --- whether it
might not be called in doubt. When a judgment is determined
more closely in one of these respects, we get a judgment about it,
hence a judgment in the second power. Such a judgment about a judgment
is called a \textit{secondary judgment}, and in contrast to it the
first judgment is called a primary judgment. Various
classes of secondary judgments have been set up, and there has been a tendency to regard
them as wholly independent kinds of judgments, which were not supposed to
be reducible to one another. But it will turn out that their
differences depend for the most part only on the linguistic
expressions, and that, when we keep strictly to their logical
significance, they can be formulated as primary judgments. It is
the same kind of judgments we have before us, whether their subject
is a single concept or a judgment.

16. Since intuition (sensation, memory, imagination) always gives
us something single and individual, a primary judgment will concern only
a single case. Thinking first discovers a partial,
absolute or limited relation of identity between the intensions of two concepts;
later it must be decided whether this relation
may be assumed to hold in several cases, or perhaps in all cases
of a similar kind. In $A = AB$, $A$ is a concept whose extension
is smaller than $B$'s (cf.\ 14 a). $A$ is more individual than $B$.
The investigation of whether $A$ can be made a typical or general
concept takes place through a series of judgments in which the intension of the subject concept
is compared with that of other, related concepts.
I must, e.g., investigate whether other human beings show
the same conduct as the man whom I (in 14 a) called good, so that
I can set up the judgment: ``some human beings are good'' [nogle],
or perhaps ``all human beings are good'' [alle]. If this widening
of the extension of the subject concept is undertaken, the determination of quantity
(some or all) which is introduced becomes a logical predicate of the judgment,
however it may be expressed grammatically. A
% --- p. 19 ---
\opage{19}\textit{quantitative} judgment is a judgment about in how many cases a certain given
judgment holds; in this secondary judgment the given judgment is the
logical subject. Grammatically the determination of quantity is added to
the linguistic subject of the given judgment (cf.\ 12: ``All the guests
have come,'' i.e., the guests now arrived make up all those who
were expected). --- Often, however, the intension of the subject concept must be diminished
before the judgment can be extended to hold for more phenomena:
the concept of this particular human being has a greater
intension than the concepts ``some'' [nogle] and ``all'' [alle] human beings.

A quantitative judgment can be expressed by limited identity.
If I have found that the judgment $A = AB$ holds only in some cases,
I can denote this by $CA = CAB$, $C$ meaning the
qualities peculiar to these cases.

There is no ground for setting up quantitative judgments as
a special kind of judgments.

17. When a judgment is put forward as unconditionally valid, it is called
\textit{categorical}; when it is put forward as conditioned by the
validity of another judgment, it is called \textit{hypothetical}.

On closer investigation it turns out that every judgment
can be formulated as a hypothetical judgment, since there can always
be pointed out presuppositions that must be present in order that
the judgment may hold. (Cf.\ 14 c.) If one wishes to formulate
these presuppositions, this can be done by setting up a limited
relation of identity between the concepts connected in the judgment:
in $AC = BC$, $C$ denotes the totality of the presuppositions that
must be present in order that $A = B$ may hold. However, a hypothetical
judgment can also be expressed in another way. The hypothetical
judgment ``If $A$ is $B$, then $C D$'' asserts that the judgment ``$C$ is
$D$'' holds only if the judgment ``$A$ is $B$'' holds. It is thus
a judgment about a judgment, a secondary judgment. The two judgments in a
hypothetical judgment are related as subject and predicate in a primary
judgment. If we now say that the conditions for $A$ being $B$ are denoted by $P$, and the conditions for
$C$ being $D$ by $Q$, then that hypothetical
judgment asserts an identity of $P$ and $Q$: either $P = PQ$ or $P = Q$ or
$PR = PQR$ (or $QR$), $R$ in the last case denoting the
more remote conditions for an identity of $P$ and $Q$.

Often the difference between a categorical and a hypothetical
% --- p. 20 ---
\opage{20}judgment is purely linguistic. The sentence: ``When the sides of a triangle
are equal, the angles are also,'' can also be expressed
thus: ``An equilateral triangle is equiangular.''

18. Immediate judgments are positive. Greater
experience (more intuitions) belongs to setting up \textit{negating} judgments than
to setting up positive judgments. A negating judgment is set up only
when there is occasion to reject a positive expectation.
When I enter a room and say: ``There is no one here!'' this
negating judgment is a rejection of an expectation that
someone was in there, and it is due to a comparison
between the room as it presents itself to my view, and the room
as I expected to find it. There prove to be fewer
elements in the observation than in the expectation. (Cf.\ on
negating concepts: 8.) A negating judgment is thus a judgment
about a judgment, a secondary judgment.

A negating judgment (e.g.\ $A$ is not $B$) can on closer
consideration be formulated in just as simple a way as a
positive judgment. For it expresses a relation of identity (partial
or absolute or limited) between the one concept and
the negation of the other concept. If we let the sign $\sim$ mean ``is not,''
and denote (as in 8) by small letters the negations
of the concepts denoted by the corresponding capital letters
(so that e.g.\ $a$ means the absence of the marks denoted by $A$),
then, if $A \sim B$, $A = Ab$ and $B = Ba$. There
is thus partial identity between $A$ and $b$, and between $B$ and $a$.
$A$ and $B$ are then said to stand in a \textit{contrary relation} to one another: they
are incompatible, they exclude one another; – but because the one
concept does not hold, the other concept need not hold.
There are more than two possibilities. What is not red need
therefore not be blue. What is not good need
therefore not be evil. A right-angled triangle is not
equilateral; but because a triangle is not equilateral, it is not
thereby said that it is right-angled --- $A \sim B$ can nevertheless sometimes mean
$A = b$. Examples of this are found wherever a
choice must be made between two possibilities. According to the strict Stoic conception
everything was evil that was not good. The relation between position
and negation (Yes and No) is of this kind; a
% --- p. 21 ---
\opage{21}question must (when one acknowledges the presuppositions under
which it is put) be answered either affirmatively or negatively.
If I negate $A$, I get $a$; if I negate $a$, I get $A$ again; I
thus do not get beyond these two possibilities. Such a
relation is called \textit{contradictory}. There is then not merely a relation of opposition,
but a relation of contradiction between the two concepts.
Not only do they exclude one another, but when the one
does not hold, the other necessarily holds. A contradictory
relation is a two-termed relation, whereas a contrary relation
is a many-termed relation.

Sometimes a relation which under ordinary circumstances
is contrary may under certain special conditions pass
over into being contradictory. We then get $AC = bC$. One
must grant the Stoics that under certain conditions one
acts evilly if one does not act well, e.g.\ if one can
save a human being and does not do so.

19. When over against a given judgment another
judgment can be set, whose validity excludes the validity of the former, there being a
contrary or contradictory relation between the two judgments,
--- then the given judgment becomes \textit{problematic} and makes the
other judgment problematic too. The two judgments which Galilei discusses in his
``Dialoger'': ``The earth moves'' – ``The earth does not move,''
made one another problematic. In this example the relation
between the two judgments is contradictory. On the other hand we have a
contrary relation when the two judgments: ``The table is black'' and
``The table is brown'' are set against each other; for here there are more
than two possibilities.

Problematic judgments are secondary. They are properly a
kind of hypothetical judgments. A judgment is problematic when its
validity presupposes the invalidity of another judgment, and it can
therefore be formulated like a hypothetical judgment with a negative presupposition
(thus $P = Pq$ etc., see 17). Yet in a problematic
judgment the two judgments stand more side by side than in the
hypothetical, where the relation is not always reciprocal (except in
$P = Q$). A problem arises when two mutually exclusive
judgments are possible, so that one stands before an Either---Or.

If proof can be given for one of the judgments, the other falls
% --- p. 22 ---
\opage{22}away, and the problematic character is removed. The same may
be attained when one of the judgments is refuted. But if the relation
is contrary (not contradictory), there may be several judgments
that conflict with the given one, and which therefore must also be
got out of the way, if the latter is to lose its problematic character.
--- A judgment that has lost its problematic character in one of these ways
is called apodictic.

20. Common to all secondary judgments (the quantitative,
hypothetical, negative and problematic) is that the auxiliary
synthesis, which in a primary judgment is only of subordinate significance
(10 and 14 a), plays a more prominent role.
When I am to determine the extension in which a judgment holds
(16), and the presuppositions under which it holds (17), I must
expressly hold together and compare it with other
intuitions or judgments. And the same is necessary when a
judgment is to be removed (18), and when its relation of opposition to
another judgment is to be demonstrated (19). --- Yet it is always merely
a difference of degree that there is between primary and secondary judgments.
This is seen from the fact that they can be formulated in quite the same way.

% --- p. 23 ---
\opage{23}\chaphead{III}{ANALYSIS OF JUDGMENTS}{FORMATION OF INFERENCES}


\sechead{A. The Logical Principles.}{\textbf{A. The Logical Principles.}}


21. Already in the transition from intuition to concept and
judgment, a principle asserts itself, whether one is conscious of it or not,
namely that the intension of the concept or
judgment is to correspond to the intension of the intuition, so that
this transition can be regarded as a conversion into another
form. In the formation of concepts and judgments nothing new is produced in a
certain sense, except precisely the new form itself, which
becomes significant for the continued work of the mind. But
the same principle appears still more clearly
in inference, i.e. in the formation of new judgments out of given
judgments. The conclusion (Konklusionen) is to correspond to the premisses
(Præmisserne) and to be only a conversion of
their intension into a new form. The relation between conclusion and
premisses is analogous to that between concept or judgment and
intuition. It is the relation called the relation between
\textit{ground and consequence}. In the formation of inferences this relation
appears distinctly; therefore the logical principles are discussed as an
introduction to the theory of inference, although they can already be traced
in the theory of concept and judgment.

22. The most general principle of thinking
is that our concepts and judgments are to be \textit{grounded}, i.e. traced
back either to given intuitions or to other concepts and
judgments that are assumed valid. Our concepts and judgments we form
at first involuntarily, and the task of grounding them comes
afterwards. The first demand of logic is that concepts and
% --- p. 24 ---
\opage{24}judgments are to have their sufficient ground: \textit{the principle of
sufficient ground} (principium rationis sufficientis). How
far we are to go back, and can go back, in grounding depends
on the nature of the particular question. Often we must stop
for lack of data; but in itself the said
principle leads into an infinite series of tasks.

Whether there is in reality a ground for everything that happens
is a question that formal logic, which has to do only
with the ordering of our thoughts, cannot decide. A distinction must
be drawn between ground and cause. Concepts and judgments have
their ground in intuitions; conclusions have their ground
in judgments that are assumed valid. But the causal relation is assumed
to obtain between events (in us or outside us)
that succeed one another in time, and which are so linked to
one another that when the one occurs, the other also occurs inevitably.
Between premisses and conclusion
there is no temporal relation; the latter is already contained in (``lies in'')
the former, even though we may be long in seeing it. But in
the causal relation the event that is cause precedes in time
the one that is effect. Between the concepts ground and
consequence on the one hand, and the concepts cause and effect on
the other hand, there is an analogy, which is expressed by saying that the cause
is called the ground of becoming of the effect (besides being for
us a ground of knowledge, i.e. a ground for expecting the event
we call the effect). The problem of cause, which it is
the task of the theory of knowledge, not of formal logic, to discuss,
concerns precisely the justification of the analogy between the grounding of our propositions
and the temporal succession of events. --- In formal
logic and pure mathematics we have to do only with grounds;
in the real sciences (natural science, psychology and history)
we have to do with causes as well. The difference between the
formal and the real sciences depends on the role that differences of quality
and differences of time play in the latter, whereas the formal
sciences abstract from them. (For formal logic,
however, figures and magnitudes are also qualities, cf. 2.)

23. Our striving for grounding can stop at two
points, either at the intuitions from which we form the
% --- p. 25 ---
\opage{25}simplest concepts and judgments, or at the principles
that are presupposed in the formation of concepts and judgments and
in the drawing of inferences. In the first case it must be examined
whether our sense-intuitions are not illusory, whether our
memories do not deceive us, and whether we can and will hold fast to
our images of the imagination. In the second case a continued discussion
will be impossible if any of the participants does not acknowledge
the logical principles; these, since they are presupposed
in all grounding, cannot themselves be grounded.
He who would deny the presuppositions of all grounding would place himself outside
all discussion. We discover the logical principles by
noting the procedure that is used in deriving
judgments from other judgments or from intuitions. He who
denies them must either point to another procedure
or give up all discussion.

In the grounding of judgments and the setting up of concepts
five principles may be necessary.

24. \textsuperscript{1}) The principle of identity: $A = A$. This principle is already implicit
in the definition of a concept as a fixed and universally valid
idea, i.e. an idea that, wherever it occurs,
has the same intension, and thus can be set equal to itself.
It is implicit furthermore in the relation between subject and predicate
in the judgment, when this is formulated as identity. And it is presupposed
in all inferring: from two judgments ($A = B$ and $B = C$) I
can derive a third ($A = C$) only when they have a common concept;
I can then put the expression for this concept from the
one judgment into the other. The inference takes place by substitution,
and identity can be defined as possibility of substitution. Identical
concepts must be capable of being put everywhere in the place of one another. I
have the right to put $A = B$ (or $A = AB$ etc.) only when the
concepts that are connected by the sign of equality are merely different
expressions for one and the same intension, so that I can everywhere put
the one in the place of the other.

Our actual ideas can never be shown
to be absolutely identical. We can only point to an actual
similarity of coincidence, where no difference is noticed (cf.
Psykologi VB, 8. 11); but absolute certainty of identity
% --- p. 26 ---
\opage{26}cannot be attained. Even where we think one and the same thought again
(e.g. pass from $A$ on the one side of the sign of equality to $A$ on
the other side of the sign of equality), there may be a difference of nuance
in the manner in which we think it. In setting up
the principle of identity, logic gives us a standard or an ideal,
by which we can test the exactness and validity of our concepts, judgments and inferences.
Only an approximation in the direction
of this ideal can be possible.

25. \textsuperscript{2}) \textit{The principle of simplicity}: $AAA ..... = A$. The repetition of a concept
and its combination with itself has no logical significance,
since its intension is not changed thereby. --- Whereas logic
can get nothing out of the combination of identical concepts,
it is otherwise in arithmetic: here weight is laid on the fact
that the repeated concepts presuppose just as many different
acts of thought, and that the concept of number therefore finds application just as well
in the series $AAA$ .... as in the series $ABC$ .... In the
application of the concept of number it may be a matter of indifference whether the units have
the same or different intension. --- For psychology the repetition of an element of consciousness
is of great significance; it makes recognition
possible; it tends to weaken feeling, but
on the other hand can itself be an expression of strong feeling,
as when we e.g. repeat the words in order that feeling may find vent
(Verily, verily I say unto you!). (Cf. Psykologi VI E,
4; F, 4.)

26. \textsuperscript{3}) \textit{The principle of contradiction} (the law of contradiction):
$Aa = 0$. I cannot form a combination or a judgment
by connecting a concept with its own negation. Where such
a connection (e.g. ``a round square'' or ``the square is
round'') has been formed, it must immediately be canceled. For logic
a self-contradictory connection is nothing (= 0). --- For psychology
actual self-contradictions are phenomena that are to be
explained. Self-contradiction is possible so long as the contradicting
ideas are not apprehended definitely, not analyzed exactly,
and not brought into immediate connection with one another. Many kinds of
psychological causes (especially the influence of feeling, see
Psykologi VI F, 1) can prevent consciousness from discovering
the self-contradiction.

% --- p. 27 ---
\opage{27}27. \textsuperscript{4}) While I cannot combine $A$ with $a$, I can
combine it with another concept. And I then have, as regards
every such concept, the choice between two possibilities:
to connect $A$ with this concept (e.g. $B$) or not to connect
$A$ with it. There are only these two possibilities; we thus face
here a contradictory relation. There are only two logical
dimensions: position and negation. (Cf. 18.) This
relation is expressed in the \textit{principle of duality} (principium exclusi
medii inter duo contradictoria): The concept $A$ must either be connected
with $B$ or not be connected with $B$. This can be expressed thus:
$AB \udel Ab$ (the sign $\udel$ signifying a mutual relation of exclusion,
a contradictory relation, an Either---Or).
The principle of duality expresses both the incompatibility of the two
combinations and the fact that the cancellation of the one is
the assertion of the other.

A vertebrate must either be a mammal or not be
one: the marks that are peculiar to the vertebrate type
must either occur together with those that are peculiar
to the mammal type, or without these.

Every concept and its negation divide the world into two parts
($B$ and $b$). What does not have the marks of the one concept
must have those of the other. And the same thing is repeated within
each single concept, when a division is to take place. The
negative concepts therefore play an important role in attempts
at classification. (Cf. 8.)

When $B$ is not a necessary determination of $A$, both
$AB$ and $Ab$ can occur --- under different conditions or
at different times. The vertebrate type occurs both in
mammals and in non-mammals, e.g. both in the horse and
in the eagle.

28. \textsuperscript{5}) \textit{The principle of commutativity}: $AB = BA$. (Cf. 14 b.)
It is a matter of indifference whether I begin with the one or with the
other concept in a combination, for the intension is not changed
thereby. The order of the members is a matter of indifference. --- The same rule holds
in arithmetic as regards addition and multiplication.
On the other hand, in the psychological respect (as everywhere where
real events are in question) it is not a matter of indifference in which
% --- p. 28 ---
\opage{28}order the members follow one another, since the preceding members
come to influence the following ones. (Psykologi VA, 4; VI E.)

29. The logical principles are not merely presuppositions for
all grounding, and thereby for all discussion, but they express
also the nature of human thinking. To think is to connect
through comparison (see 1 and Psykologi VB, 11), and
thinking moves forward through identifications, substitutions,
negations and combinations, the basic rules of which are
formulated in the logical principles. Finally, the logical
principles lead to the setting up of ideals of knowledge and act as
guiding thoughts in our investigations, in that they lead us to
apprehend both similarities and differences as exactly as possible
and thereby make our connections of thought ever more secure and more perfect.
Whether they hold for existence as a whole, so that this,
seen from a human standpoint, is rational, is a question
that belongs partly to the theory of knowledge, partly to cosmology
(metaphysics). (Cf. the section ``De filosofiske Problemer''
in ``Kort Oversigt over den nyere Filosofis Historie''.)

\sechead{B. Inferences.}{\textbf{B. Inferences.}}


30. An inference aims at bringing to distinct consciousness
what is contained in certain given judgments. Thereby
no new intension is produced, but we become conscious of what
we know. Sometimes we know more, sometimes less, than we thought we
knew. But it can mean a very great advance that we
become conscious of the significance and range of our judgments.

Inference can be of two kinds. In \textit{direct inference}
two or more judgments are given, and the task is then to find
a common concept (the middle term), so that substitution
can take place from the one judgment to the other. Here one proceeds
thus, by means of the principle of identity, straight from premisses
to conclusion. In \textit{indirect inference} it is investigated
which combinations of the concepts contained in one or more judgments
can subsist with the validity of these judgments.
This investigation takes place by rejecting the combinations
that, by substitution from the given judgments, prove to be 
% --- p. 29 ---
\opage{29}self-contradictory. Here, then, one builds not merely on the principle of identity,
but also on the principle of contradiction (and, as will appear later,
the principle of duality), and the result is reached through the casting out of
the invalid combinations. --- In both kinds
of inference the principles of simplicity and of commutativity
may find application.

Inference rests, like the formation of concepts and the formation of judgments,
on the interaction of analysis and synthesis. The given judgments,
which are set up and put together by a process corresponding
to the first synthesis (1. 4. 10), must be analyzed in order that one
may find which concepts can be substituted. But next,
by means of substitution, a new synthesis takes place,
which leads to seeing which judgments or combinations
correspond to and subsist with the given judgments. Inference is, as
Aristotle already expressed it, a \textit{thinking together} of the
given. The Greek word for inference (syllogismos) means
precisely thinking together. Inference is a logical synthesis.

\sechead{a. Direct Inference.}{\textbf{a. Direct Inference.} (Inference by Substitution Alone.)}


31. It is presupposed here that we have (in the simplest case)
two judgments given that have a common concept, so that substitution
from the one judgment into the other becomes possible. Such a concept
is called the middle term; the concepts in relation to which the middle term
is either subject or predicate are called the principal concepts.
--- Of the innumerable forms of direct inference
only a few are given here as examples.

32. \textit{When a subject has a predicate that in turn is the subject
of another predicate, it is inferred that that subject also has
this other predicate}:
\[\begin{array}{l}
A = AB \\
B = BC \\
\hline
A = ABC = AC
\end{array}\]
Neptune is a planet; every planet moves in an ellipse;
therefore Neptune also moves in an ellipse. A square
is a parallelogram; in a parallelogram the
% --- p. 30 ---
\opage{30}diagonals bisect each other; therefore in a square the diagonals bisect
each other.

This inference is expressed in the Aristotelian logic, which
builds on the extension of concepts (see 13), thus:
\[\begin{array}{l}
A < B \\
B < C \\
\hline
A < C.
\end{array}\]
It is called an inference in the \textit{first figure}. Aristotle distinguishes
between different ``figures'' (schemata), in that, according to his exposition,
the different place of the concepts as subject or predicate
in the premisses became of great importance. In the first figure the
middle term is predicate in the one premiss, subject in the
other. --- In geometrical representation $A$ will be a small
circle inside $B$, and $B$ in turn a circle inside $C$, and it is then seen
immediately that $A$ is contained in $C$. ---

If instead of partial identities we had absolute identities
in the premisses, the conclusion would also be an absolute
identity:
\[\begin{array}{l}
A = B \\
B = C \\
\hline
A = C.
\end{array}\]
This form of inference (which could not occur in Aristotle
and cannot be represented geometrically) is used when the identity
of phenomena is to be established. If I have an exact description
($B$) of a man ($A$), and this description fits entirely
the person ($C$) who committed a certain murder, I infer
that $A$ is the murderer. The force ($A$) that holds the moon in
its orbit is the same force ($C$) that on earth is called
gravity, because the law ($B$) for both forces is the same. ---

If I have
\[\begin{array}{l}
A = AB \\
B = BA \\
\hline
\end{array}\]
there follows $A = B$.

% --- p. 31 ---
\opage{31}Here, besides the principle of identity, the principle of
commutativity is also applied. An absolute identity is inferred here from two partial
identities. What we get here is, properly, only a definition of absolute
identity. (Cf. 14 b and 28.) --- An equiangular triangle
is equilateral, and an equilateral triangle is equiangular; the two
kinds of triangles are therefore identical.

33. \textit{From the fact that two subjects have one and the same predicate,
nothing can be inferred about the relation between them}.
\[\begin{array}{l}
A = AB \\
C = CB
\end{array}\]
Here there is nothing to substitute. In Aristotle this becomes the
\textit{second figure}:
\[\begin{array}{l}
A < B \\
C < B,
\end{array}\]
and he warns expressly against inferring here, unless
the one premiss is negative. (When $A$ is $B$, but $C$ is not
$B$, $A$ cannot be $C$.) --- Geometrically both $A$ and $C$ lie
inside the circle $B$; but their mutual relation cannot be
decided; they may coincide or intersect or
fall outside each other or the one inside the other.
--- From the fact that both fishes and whales live in water, it can neither
be inferred that whales are fishes, nor conversely. But from the fact that fishes
have cold blood and whales warm blood, it can be inferred that whales
are not fishes. If $A = AB$ and $C = Cb$, $C$ cannot be
predicate of $A$, since we should then get $ABCb$, which contains a
contradiction.

34. \textit{When one and the same subject has two predicates,
a limited identity of the two predicate concepts can
be inferred from this}.
\[\begin{array}{l}
B = BA \\
B = BC \\
\hline
BA = BC
\end{array}\]
Fishes are both aquatic animals and vertebrates, and these two
properties then occur together under certain conditions.
% --- p. 32 ---
\opage{32}Man is a spiritual being, and man is a material
being; the spiritual and the material must then be
compatible under certain conditions. (Whether the conditions that are
expressed by $B$ are the only ones under which $A$ and $C$ can
occur together is not said here.)

This inference is represented in the Aristotelian logic thus:
\[\begin{array}{l}
B < A \\
B < C
\end{array}\]
and is called an inference in the \textit{third figure}. The middle term is here
subject in both premisses. The conclusion becomes: some $A$ is
$C$ (e.g.: some aquatic animals are vertebrates). --- Geometrically the
third figure can be represented by the circles $A$ and $C$ intersecting,
and $B$ being the piece that belongs to both circles.
(If $A < C$, we return to the first figure.) ---

If we have not merely $B = BA$ and $B = BC$, but also
$A = AB$ and $C = CB$ (and so can combine the 2nd and 3rd figures),
we get $A = C$, so that the two concepts, which at first seemed only to
be limitedly identical, now prove to be absolutely identical.
If $A$ means ``most highly developed rational being'', $B$ ``man''
and $C$ ``animal with the most perfect brain'', we get that the
most highly developed rational being is the animal that has the most perfect
brain.

35. \textit{If both premisses are negative}, a conclusion can be found
by using what is set out in 18 about
negative judgments.

If I have the two negative judgments $A \sim B$ and $B \sim C$, I must
examine whether each of them is merely contrary, or whether they
are contradictory. Each of them must necessarily
be contrary.

$A \sim B$ must mean \textsuperscript{1}) $B = Ba$, but can also mean \textsuperscript{2}) $B = a$.

$B \sim C$ must mean \textsuperscript{3}) $B = Bc$, but can also mean \textsuperscript{4}) $B = c$.

If only \textsuperscript{1}) and \textsuperscript{3}) hold, we get $Ba = Bc$. If \textsuperscript{1}) and \textsuperscript{4}) hold,
we get $Ba = c$. If \textsuperscript{2}) and \textsuperscript{3}) hold, we get $a = Bc$. If \textsuperscript{2}) and \textsuperscript{4})
hold, we get $a = c$.

Here we have started from the middle term; but one can also
start from the principal concepts:

% --- p. 33 ---
\opage{33}$A \sim B$ must mean \textsuperscript{1}) $A = Ab$, but can also mean \textsuperscript{2}) $A = b$.

$B \sim C$ must mean \textsuperscript{3}) $C = Cb$, but can also mean \textsuperscript{4}) $C = b$.

If \textsuperscript{1}) and \textsuperscript{3}) hold, nothing can be inferred (since that is the second figure).
If \textsuperscript{1}) and \textsuperscript{4}) hold, we get $A = AC$. If \textsuperscript{2}) and \textsuperscript{3}) hold,
we get $C = CA$. If \textsuperscript{2}) and \textsuperscript{4}) hold, we get $A = C$.

36. Sometimes only one premiss is expressly
given, and yet an inference can be drawn. Such an inference
is called an \textit{immediate inference}. It stands as self-evident.
But if a proof is demanded, one or more of the
logical principles are needed for it. Here (in direct inference) we have to do only with
the immediate inferences that can be proved by
the principle of identity.

If $A = B$, then $A = AB$ and $B = BA$ follow. This lies
already in the very concept of absolute identity. (Cf. 14 b. 28. 32.)
But it can be proved thus: $A = AA$ (by the principle of
simplicity). If now $B$ is substituted in one place on the right-hand side, we get $A = AB$.
In the same way we get from $B = BB$ that $B = BA$. --- From $A = B$
it can also be inferred that $AC = BC$. What holds absolutely must
also hold under a certain condition. The proof is as follows: We
put: $AC = AC$ and insert on the one side $B$ for $A$, so
that we get $AC = BC$.

37. \textit{The question whether an inference can be reversed}, so
that the conclusion becomes premiss, and the earlier premisses conclusions,
must be examined in the various cases.

The inference from $A = AB$ and $B = BA$ to $A = B$ (33) can be
reversed (see 36). But the inference from $A = B$ to $AC = BC$
 (36) cannot be reversed, since one has no right without further ado
to cancel the condition. --- Where the middle term disappears
in the conclusion, as when (32) one infers from $A = AB$ and $B = BC$
to $A = AC$, or from $A = B$ and $B = C$ to $A = C$, the inference
cannot be reversed, since one cannot find the middle term from the
conclusion, and since, moreover, any middle term whatever
will give the same conclusion. $A = C$ follows just as well
from $A = X$ and $X = C$ as from $A = B$ and $B = C$.

The ground (the premisses) leads with necessity to
the consequence (the conclusion), but cannot always go the opposite
way. In empirical science one infers back from a given
% --- p. 34 ---
\opage{34}phenomenon (e.g. the presence of carbonic acid) to the conditions for its
origination; but these may be different in different cases.
(Carbonic acid may arise through the respiration of animals, through
combustion, through volcanic eruptions etc.) Therefore most
explanations of nature are hypotheses, grounded conjectures.
And a hypothesis does not become fully certain because it is confirmed
by experience, so long as it cannot be shown that
no other premisses could lead to the same result,
--- so long, that is, as the inference cannot be reversed. Leibniz
already saw this: le succès ne démontre pas la vérité
de l'hypothèse. All modern natural science thus rests
on presuppositions that were set up by Galileo and his
successors, and the question is whether these presuppositions are
the only ones from which the results won can be reached.
(Cf. on this the more recent viewpoints in the theory of knowledge mentioned in ``Oversigt over den nyere Filosofis
Historie'' VIII, 2.)

Not only can one and the same conclusion be derived from judgments
with highly different middle terms, but a correct conclusion
can be derived from factually incorrect or even self-contradictory
premisses. Both $A = B$ and $B = C$ may be factually incorrect,
and yet $A = C$ may be factually correct. It often turns out
that one arrives from false premisses at a result that
is then proved correct by other routes. The correct derivation
of $A = C$ might e.g. be that from $A = K$ and $C = K$. ---
From $A = a$ and $a = C$ follows $A = C$, which contains no
self-contradiction, although the one premiss is self-contradictory.

\sechead{b. Indirect Inference.}{\textbf{b. Indirect Inference.} (Inference by Substitution and Elimination.)}

38. Indirect inference investigates which combinations of concepts
can subsist with the validity of one or more
given judgments. This investigation takes place by inserting, in the
combinations (according to the principle of duality) conceivable of the concepts given in the judgment
(or judgments), the expressions
that the judgments give for these concepts, and casting out (eliminating)
the combinations that become self-contradictory. This is a deduction,
in that one goes from one or more given judgments to the
combinations that can subsist with them. --- The judgments
% --- p. 35 ---
\opage{35}that are investigated in this way may themselves have been derived
by \textit{induction} from certain given combinations of concepts. Deduction
then provides the proof that the judgments have been correctly derived.
Induction is essentially a method of discovery; by it we procure
judgments that can be tested deductively. It is preparatory
work; the inference proper is always deductive.

\sechead{α. Indirect Inference as Pure Deduction.}{\spaced{α. Indirect Inference as Pure Deduction.}}


39. When the judgment $A = B$ is given, which combinations
of $A$ and $B$ and their negations will then be possible? --- Since
two concepts are in question, only four combinations can be conceived
\textit{according to the principle of duality}: $A$ can occur either with
$B$ or without $B$; $B$ either with $A$ or without $A$; $a$ either with
$B$ or without $B$; $b$ either with $A$ or without $A$. The four combinations
that can be conceived are thus $AB$, $Ab$, $aB$ and $ab$.
(Geometrically this can be illustrated by letting the circles
$A$ and $B$ intersect; there is then a part of $A$ without $B$,
a part that belongs to both $A$ and $B$, and a part of $B$ without $A$;
to this is added the space outside the two circles, where neither
$A$ nor $B$ is found.) But now the task is to find which of
these conceivable combinations are possible when the judgment
$A = B$ is given. This is done by inserting this judgment in the
four combinations \textit{(by virtue of the principle of identity)} and then
\textit{(by virtue of the principle of contradiction)} casting out the combinations
that become self-contradictory. The combinations in which
nothing can be substituted remain standing as possible.
\[\begin{array}{l}
AB = AA = BB \\
Ab = Bb = 0 \\
aB = aA = 0 \\
ab
\end{array}\]
Thus only $AB$ and $ab$ can occur when $A = B$. ---
If the person $A$ is to be identical with the murderer $B$, then
he must have been present everywhere the murderer has been,
and have been absent where the murderer was absent.

40. When the judgment $A = AB$ is given, we get by the same
procedure $AB$, $aB$ and $ab$ as possible cases. --- If
% --- p. 36 ---
\opage{36}heat is a form of motion, we must be able to find phenomena of heat
and motion together, or motion without heat
(since there are other forms of motion than heat), or phenomena
that are neither motion nor heat. Only the combination
heat without motion ($Ab$) is excluded.

When $A = Ab$ is given, $Ab$, $aB$ and $ab$ become possible. When
$A = b$ is given, only $Ab$ and $aB$ become possible.

41. For simplicity's sake examples with two concepts have been chosen in the foregoing.
If three concepts are in question,
we get 8 conceivable combinations to examine. (Geometrically
this can be represented by letting three circles intersect
one another so that each of them intersects both the others;
one then gets 7 different pieces besides the space outside
the three circles.) If four concepts are contained in the given
judgments, we get 16 combinations etc. --- The following example
with three concepts shows the procedure.

Let $A = AB$ and $C = CB$ be given. Then the 8 conceivable
combinations are:
\begin{center}
\begin{tabular}{l@{\qquad}l}
\textsuperscript{1}) $ABC$ & \textsuperscript{5}) $aBC$ \\
\textsuperscript{2}) $ABc$ & \textsuperscript{6}) $aBc$ \\
\textsuperscript{3}) $AbC$ & \textsuperscript{7}) $abC$ \\
\textsuperscript{4}) $Abc$ & \textsuperscript{8}) $abc$
\end{tabular}
\par\smallskip By substitution we get:\par\smallskip
\begin{tabular}{l@{\qquad}l}
\textsuperscript{1}) $ABBCB = ABC$ & \textsuperscript{5}) $aBCB = aBC$ \\
\textsuperscript{2}) $ABBc = ABc$ & \textsuperscript{6}) $aBc$ \\
\textsuperscript{3}) $ABbCB = 0$ & \textsuperscript{7}) $abCB = 0$ \\
\textsuperscript{4}) $ABbc = 0$ & \textsuperscript{8}) $abc$
\end{tabular}
\end{center}
By indirect inference one can thus draw inferences from the
``second figure'' (33) concerning the combinations of concepts that
are possible when two subjects have the same predicate.

If one starts from $B = BA$ and $B = BC$ (thus the ``third figure'',
see 34), $ABC$, $AbC$, $Abc$, $abC$ and $abc$ become possible.

42. If we start from mutually contradictory premisses,
indirect inference can call attention to the error.

% --- p. 37 ---
\opage{37}
\[\begin{array}{l@{\qquad}l}
A = b \\
B = A \\
\hline
AB = bB = 0 & aB = aA = 0 \\
Ab = bB = 0 & ab = aA = 0
\end{array}\]
All conceivable combinations thus fall away. --- Another
example will occur below (46).

\whiteline

\sechead{β. Indirect Inference as a Union of Induction and Deduction.}{\spaced{β. Indirect Inference as a Union of Induction}\\
\spaced{and Deduction.}}

43. If we have a series of combinations of concepts given, we can
seek one or more judgments that express the relation between
the concepts contained in the combinations. Such judgments
are found by analysis, as when we form our judgments on the
basis of immediate intuition (10). When the combinations
are very numerous, one must feel one's way forward through provisional
hypotheses. Thus Kepler first tried the assumption that
the orbit of Mars was a circle, before he discovered that it was an
ellipse. The auxiliary synthesis plays a great role here.
The combinations stand here side by side, whereas in the simpler
cases of the formation of judgments the task is to illuminate a single
combination (given in intuition) by means of the others.
(Cf. 10 and 20.)

The transition from given combinations to judgments is called
induction. But a proof that the judgments found really
correspond to the given combinations and not to others must be given
by deduction. The proof of an induction is always a deduction,
and the significance of induction consists in setting up a conjecture
that can be tested. (Cf. 38.)

Indirect inferences must always be capable of being reversed, if they
are correct. I must be able to go both from a definite judgment to
definite combinations, and conversely.

44. If the combinations $AB$ and $ab$ are given, we see that where
$A$ is, $B$ is, and conversely, and that where one of them is lacking,
the other is lacking too. To this correspond $A = B$ and $a = b$. And
by deduction from these judgments we come also to $AB$ and
$ab$, and to them alone (see 39).

% --- p. 38 ---
\opage{38}If $AB$, $aB$ and $ab$ are given, we see that where $A$ is, $B$ is, whereas
$B$ can occur without $A$, and that where $B$ is lacking, $A$ is lacking.
That gives $A = AB$ and $b = ba$. Of $B$ and $a$ nothing can be asserted.
Deduction from $A = AB$ leads also to $AB$, $aB$ and $ab$, whereas
$Ab$ would be self-contradictory. (See 40.)

If $AB$, $Ab$ and $aB$ are given, it is seen that $a$ occurs only if
$B$ is, whereas $B$ can occur without $a$, and that where $B$ is lacking,
$a$ is lacking too. To this correspond $a = aB$ and $b = bA$. Of $A$ and
$B$ nothing can be asserted here. By deduction from these two judgments
we get also the three given combinations and no others.

If the given combinations were $ABC$, $AbC$, $aBc$ and
$abC$, there would lie in them $A = AC$, $b = bC$, $c = cB$, furthermore
$AB = ABC$ etc.

45. Many \textit{immediate inferences} (cf. 36) can be
proved only by indirect inference. This holds of the important kind
of immediate inference by which \textit{from a positive judgment we
derive a negative judgment or conversely}. Such a derivation
is called \textit{contraposition}.

From $A = AB$ follows immediately $b = ba$. For when $B$ is where
$A$ is, $A$ must be lacking where $B$ is lacking. A proof can nevertheless be given
thus: Of the four possible combinations, when $A = AB$
is given, $Ab$ drops out; only $AB$, $aB$ and $ab$ remain, and from
them follows (see 44) $b = ba$.

From $A = B$ follows immediately $a = b$. If there is absolute identity
between $A$ and $B$, $a$ must be lacking where $b$ is lacking, and conversely.
A proof can be given thus: Of the four possible combinations,
when $A = B$ is given, $Ab$ and $aB$ drop out; only $AB$ and $ab$
remain, and from them follows (see 44) $a = b$.

These inferences can be reversed. From $b = ba$ follows by
contraposition $A = AB$, from $a = b$ follows $A = B$. ---

By means of contraposition we can formulate an inference
in the ``second figure'' with a negative premiss (see 33):

\begin{center}$A = AB$\end{center}


\begin{center}$C = Cb$\end{center}


From $C = Cb$ follows $B = Bc$, so that we transform the second figure
into the first figure:

% --- p. 39 ---
\opage{39}
\[\begin{array}{l}
A = AB \\
B = Bc \\
\hline
A = ABc = Ac.
\end{array}\]
Even from the second figure with two positive premisses we can
now get a result. For from $A = AB$ follows $b = ba$, and from $C = CB$
follows $b = bc$, so that we get:
\[\begin{array}{l}
b = ba \\
b = bc \\
\hline
ba = bc
\end{array}\]
The second figure is thus changed by contraposition into the third
figure. ---

46. In every induction we must have all the concepts contained
in the combinations given both in negative and in positive
form, since every concept must be capable of being negated. (Cf. 5---6.)
If we are to infer from two combinations alone, these must therefore
be either $AB$ and $ab$ or $Ab$ and $aB$ (when only two concepts
are in question).

If, on the other hand, we had e.g. $AB$ and $aB$ given, we should
get the two judgments $A = AB$ and $a = aB$. But from them follows by
contraposition (45) $b = ba$ and $b = bA$, which conflict with one another
(make one another problematic, see 19). A contradiction thus
arises when we infer from given combinations that
exclude the negation of a concept. This contradiction can only
be removed when we can find more combinations. If we find
e.g. $ab$, the judgment $b = bA$ drops out, and the contradiction ceases.
Thus the problems that experience poses lead us to extend
our experience. --- Conversely, by proceeding deductively
from the two propositions $A = AB$ and $a = aB$ we should arrive at the result
that the concept $B$ could not be negated, and this result
would naturally lead us to discover the contradiction between
the premisses. ---

If all conceivable combinations are given, no
judgment can be set up. All concepts then occur in all the connections
conceivable by virtue of the principle of duality: $A$ now with $B$,
now with $b$, $B$ now with $A$, now with $a$ etc. If our
% --- p. 40 ---
\opage{40}experience were always of this nature, we could form no judgments,
and if the whole of nature were of this nature,
the causal principle would not hold, since it presupposes that a
phenomenon occurs only in connection with a certain definite
other phenomenon. Where experience seems to present cases
that offer all conceivable combinations (with two concepts
thus both $AB$, $Ab$, $aB$ and $ab$), we assume --- if we
are concerned to hold fast to the causal principle --- that our experience
is incomplete, in that it is limited to two concepts.
If our experience is extended so that we take account of a
third concept $C$, which occurs together with $A$ and $B$, we shall
perhaps be able to set up one or more judgments. Thus if
we found the combinations $ABC$, $AbC$, $aBc$ and $abC$ (see the end of 44).
If, on the other hand, all 8 combinations of $A$, $B$ and $C$
occurred, no judgment could be set up here either. We
should then have to set up a fourth concept $D$, which could belong together
with the three concepts. And so on, until the difficulty
is removed. (Cf. 14. A limited identity cannot be
derived from the subject concept and the predicate concept alone, but is found
by an extension, which establishes that the relation between the two concepts
can be fixed only when a third concept is taken in.)
--- Thus the thinking elaboration of the given leads
to seeing its incompleteness and thereby drives to extending experience.


47. Even if the judgment found leads by deduction to the
given combinations and to no others, and thus proves to
be correctly derived, it nevertheless remains a \textit{hypothesis},
so long as it has not been established with certainty that the given
combinations are exhaustive, so that no combinations
of another kind occur. Such certainty can, apart from
very simple and surveyable relations, be very difficult
to attain. The scientific character of the hypothesis depends on its
being the logical expression for all given combinations. It then stands
as a \textit{grounded} conjecture. But it needs \textit{confirmation}
(verification), so long as there is still the possibility of
new combinations. The confirmation takes place when the new combinations
are shown to be of the same kind as those on which
% --- p. 41 ---
\opage{41}one relied in setting up the hypothesis. (Cf. 5 on
definition as a hypothesis.) The hypothesis $A = B$ is confirmed
if all new experiences of $A$ and $B$ are of the kinds $AB$
or $ab$; it is refuted if the presence of $Ab$ or
$aB$ can be shown. (See also 37.)

The scientific induction starts from certain cases
as typical and shows that all other cases belong to the same
type. Newton's theory of gravitation was grounded by an investigation
of the motion of the moon; it was confirmed by the fact that Kepler's
laws could be derived from it, --- that ebb and flood could be
explained by it, --- and later by a multitude of new discoveries,
e.g. the discovery of Neptune, which was brought about by
inference from the theory of gravitation. The working value of a hypothesis
(its significance as a principle for research) coincides
with the possibility of its verification.

% --- p. 42 ---
\opage{42}\chaphead{}{APPENDIX}{}


\sechead{1. Mill's Inductive Methods.}{\textbf{1. Mill's Inductive Methods.}}


JEVONS was the first to show with clarity that the proof of an induction
consists in a deduction which, starting from provisionally set up judgments,
leads back to the combinations from which these judgments were derived,
together with a continued demonstration that all
combinations are of the same kind as those from which the judgments
were derived. ARISTOTLE understood by induction a complete
collection of a certain kind of phenomena together with the showing of judgments
that hold for them all. But such a complete collection
is very seldom possible, and the problem is precisely how our
knowledge can be built up and carried further, although it cannot be proved that the given
is exhaustive. BACON and STUART
MILL have given special descriptions of the methods of induction,
without having seen their close connection with deduction,
which alone can give a proof. Stuart Mill's exposition,
which is built on a study of the history of modern empirical
science, nevertheless has lasting interest through the clarity with which
it brings out the essential forms.

a. \textit{The method of agreement}. When a phenomenon $B$ occurs
under different circumstances, such that there is only
one thing ($A$) that always precedes, while everything else varies,
it is conjectured that $A$ is the cause of $B$ (or that both are successive effects
of the same cause). This conjecture will be the
better grounded, the more cases we have of this kind.
\[\begin{array}{l@{\qquad}l@{\qquad}l}
AVX & AYZ & ......... \\
BV_1X_1 & BY_1Z_1
\end{array}\]
% --- p. 43 ---
\opage{43}If the different cases in which the formation of dew
occurs are compared (on surfaces in the open air after sunset, on
a glass of fresh water in a warm room, on the inside of
a window in a sudden downpour etc.), they present, despite
all differences, the common feature that, before the formation of dew,
a sudden cooling of a surface in relation to the surrounding
air occurs, and in this we then find the cause
of the formation of dew.

This method can, however, lead only to a provisional conjecture.
It draws attention to the relation between
$A$ and $B$, but does not exclude the possibility that $B$ might some time
occur without $A$. (We have only the combination $AB$ given
and know nothing of the negative cases. Cf. 46.) Does condensation
not take place when sudden cooling is excluded? This
question must be answered before the inference is secure. The use
of a medicine may have proved beneficial in a series of cases;
but is it beneficial in all cases, and --- would the patients
not have been cured even if they had not
used the medicine? --- We are thereby led over to the next
method.

b. \textit{The method of difference}. When a phenomenon $B$ does not occur
unless the phenomenon $A$ precedes, while all other circumstances
are alike (so that the occurrence of $A$ and $B$ constitutes the only
difference), it is inferred that $A$ is the cause of $B$ (or that $A$ and $B$ are
effects of the same cause).
\[\begin{array}{l@{\qquad}l}
AVX & VX \\
BV_1X_1 & V_1X_1
\end{array}\]
Here only two cases are needed, one where $A$ and $B$ occur,
and another where under otherwise equal circumstances they do not
occur. --- With a clouded sky or in wind no formation of dew takes place,
only in calm and clear weather. By stretching a
sheet at a height of some feet above the ground one prevents the radiation of heat
from it, and then the formation of dew also drops away
under the sheet, while it goes on just outside. --- The task is
to find a single circumstance that is decisive
% --- p. 44 ---
\opage{44}(in that it is either the whole cause, or contains the cause,
or is a part of the cause, which always consists in a whole group
of conditions, cf. Psykologi VD, 5 d; --- this remark
holds for all four methods, but is of particular importance here).
Gray discovered that the difference between electrical conductors and
insulators depended on the constitution of the substances; he had
at first assumed that when a silk cord was a better insulator than
a hemp cord, it was because it was thinner; but it turned
out that a metal wire that was as thin as the silk cord
conducted the electricity. --- A man's blameless conduct may
be due to his good character --- but perhaps also to his fortunate
circumstances, which bring no temptations. --- After the Danish
agrarian reforms at the end of the 18th century there occurred
a great rise in the prices of agricultural products, so that one could not
see clearly what was due to them; their fortunate influence
became clear, on the other hand, when prices in the 19th century
began to fall again. ---

In the method of difference we have the combinations $AB$ and $ab$
given, from which follows $A = B$ (which is here to be interpreted thus:
if $A$ is, $B$ is, and $B$ does not occur without $A$). (Cf. 17.)

The last two methods are more special forms of the method of difference.


c. \textit{The method of residues}. When $AVX$ precedes $BV_1 X_1$, and when
it is known that $V$ is the cause of $V_1$ and $X$ of $X_1$, it is inferred that $A$ is the
cause of $B$.
\[\begin{array}{l@{\qquad}l@{\qquad}l}
AVX & V & X \\
\cdashline{1-3}
BV_1X_1 & V_1 & X_1
\end{array}\]
From the irregularities in the orbit of the planet Uranus, which
could not be explained by the attraction of the known planets,
it was inferred that beyond Uranus there must be a hitherto
unknown planet, and this led to the discovery of Neptune. --- Human
characters and racial peculiarities cannot be explained merely by
the influence of upbringing and other external (physical and social)
circumstances; there remains an unexplained residue, which
must be explained by heredity. (Cf. Psykologi VII C, 2.)

d. \textit{The method of proportional variations}. When a
% --- p. 45 ---
\opage{45}phenomenon $B$ increases or decreases, or in general changes, according
as another phenomenon $A$ undergoes corresponding changes,
while all other circumstances remain unchanged,
it is inferred that $A$ is the cause of $B$ (or that both are
effects of the same cause).
% PRINTED AS IS: the middle cell of the bottom row is A₂V₁X₁ (B₂V₁X₁ expected?), and B₁V₁X has no subscript on X (z-06-04)
\[\begin{array}{l@{\qquad}l@{\qquad}l}
A_1VX & A_2VX & A_3VX ......... \\
\cdashline{1-3}
B_1V_1X & A_2V_1X_1 & B_3V_1X_1
\end{array}\]
This method finds application where one cannot wholly
exclude a phenomenon. One then examines the cases where it
varies, and seeks to find other phenomena that under the same
circumstances vary in a corresponding manner. The method requires
that the varying series be continued as far as possible.

The law of inertia was grounded by Galileo by this method:
the smoother he made his balls and the grooves in which he let them
roll, the longer the balls went on rolling. --- That mental life
is particularly bound up with the nervous system (whether the one
is cause of the other, or the cause lies further back),
is inferred from the fact that they vary in a corresponding manner. (Psykologi
II, 6---7.) --- Statistics shows that the number of marriages rises when
grain prices fall, and falls when these rise. ---

With regard to these methods of Mill it must be noted that they
presuppose a distinct difference between cause and effect. Such a
distinct difference exists only for the provisional conception;
when research penetrates further into the relation,
it finds a continuous transition from what is called ``cause''
to what is called ``effect''. (Cf. Psykologi VD, 3.)

The methods of Mill presuppose the validity of the causal principle
and lie to that extent outside formal logic (cf. 22).
Mill conceded this and thereby came to move in a
circle, since he grounded the causal principle itself by induction
(the method of agreement). (Cf. ``Oversigt over den
nyere Filosofis Historie'' VII b.)

The question of the validity of the causal principle leads over
from formal logic to the theory of knowledge, which asks with
what right we ascribe to the principle of ground (and the logical principles contained
% --- p. 46 ---
\opage{46}in it, see 21 to 29) validity for real
existence, not merely for the ordering of our thoughts. --- In psychology, too,
we come to discuss the causal relation, namely
in the question how we distinguish between reality
and mere idea (illusion or dream). (Cf. on this Psykologi
VD, 1---2.)

\sechead{2. Analogy and Inference by Analogy.}{\textbf{2. Analogy and Inference by Analogy.}}


The innermost nerve of all logic is identity. It is by virtue of
it that we fix our ideas to concepts, tie
concepts together in judgments and go from judgments to judgments
through inference. The aim of our thinking is therefore, as far as
possible, to reduce differences to identities --- or at least
to apprehend them as if they were identities, and see how far
we can carry through this apprehension. But thought carries on
a constant struggle with the given in order to attain this aim. And
at several points analogy (similarity of relations) must take the place
of identity.

a. Our concepts rest for the most part, namely apart from
the ``mark-concepts'', on analogy between the phenomena for
which they hold (5). Only approximately can there be shown among the marks
of different phenomena given in intuition
so strict an identity as the principle of identity demands
(24). There appears here an irrational relation between the matter
and the form of knowledge, between the given and its expression in
thought. Once we have formed a concept, we must work
with it as identical with itself, until we perhaps, from increased
and clearer matter of knowledge, can form a more perfect
concept. But it must not be forgotten that it often rests from the start
on analogy, and that our definitions are often hypotheses.

b. When we now connect the concepts in judgments, and the judgments
in inferences, the connections thus established
between our thoughts cannot without further ado be assumed to reproduce the
connections that obtain between the given phenomena.
The relation between ground and consequence on the one hand and between
cause and effect on the other we have no right
to identify. We can apprehend the
% --- p. 47 ---
\opage{47}sequence of events as if it reproduced a relation between ground and consequence;
but only an analogy is it justified to assume here (22). In a
similar way matters stand with the application of mathematical
concepts and judgments to phenomena. The number series is the
simplest and most surveyable ordering of elements we know,
and we therefore naturally seek to order the phenomena in analogous
series, so that from the relations that hold between the numbers
we may infer relations between the phenomena. The so-called
mechanical conception of nature rests on the great analogy between
motions in space and series of qualitative changes.
Thus physics regards the series of colors as if it were merely
a series of quantitative changes (in the refraction of waves of ether
etc.), and the qualitative changes of substances are regarded in chemistry
as if they were merely rearrangements of unchanged
small parts. The application of the formal sciences to the
real sciences (cf. 2 and 22) is thus due to analogy.

c. In the examples just mentioned, inference by analogy appears
in its simplest form, as the method of proportional
variations. Inference from the one series (e.g. the quantitative)
can here always find its confirmation by an observation of new
members in the other series (the qualitative).

But it is not always that inference by analogy can assume this
form. Often it is an inference in the second figure (33): when two
subjects have a common property, it is natural to ask
whether they do not agree in more properties, especially when
the common property is of essential, deeply penetrating significance.
That both fishes and whales are aquatic animals may occasion
an investigation of what other properties they have in common.
While the second figure cannot ground any logical inference,
it can give the impulse to a discovery, which can then find
its confirmation by quite other routes. And when a subject has
two predicates (third figure, see 34), it is natural to ask
how these two predicates (e.g. aquatic animals and vertebrates
of the fishes) mutually determine one another; that may lead to
seeking out several other cases in which the two predicates occur
together. (One then properly makes the two predicates into
subjects, so that the third figure passes over into the second figure, which is
% --- p. 48 ---
\opage{48}the proper scheme for inference by analogy.) Thus analogy can
act as a driving force, and in a more positive way
than self-contradiction (see 46). It is a significant method of discovery.
But when the proof is to be given, it must be given
by means of identities. Newton's discovery of gravity as
a law of the world is said to have been made by means of a bold
analogy (10) between the motion of the moon and the motion of falling
on earth; but the proof was given according to the formula $A = B = C$ (32),
where $B$ signifies the identical law of the two motions.

Aptly Aristotle, who uses the word analogy only
in the sense of mathematical proportion, has called what we term
inference by analogy an inference from an example (paradeigma). We
infer from the relations in one case to the relations in another
case. With him it forms an opposition both to induction
and to deduction. --- It approaches induction the more
different examples it relies on, so that a judgment can be
grounded and perhaps later confirmed.

d. It is the great question in all analogizing whether
we have the right to transfer what holds in one respect
or for certain phenomena to other respects or other
phenomena, because a single point of similarity can be shown. In the end
the question poses itself, and in its most problematic form,
when we stand at the boundary of our knowledge and ask whether
we can from what our experience and thinking teach us about existence
draw inferences to its whole nature. Here mankind
--- in its religions and in its speculations
--- has worked with great and bold analogies, which cannot be confirmed
by experience, and whose worth also need not
rest on that.

\end{document}
